% ═══════════════════════════════════════════════════════════════
%  PART XIII — APRIL 8, 2026: THE ARITHMETIC SYNTHESIS
% ═══════════════════════════════════════════════════════════════
\clearpage
\part{The Arithmetic Synthesis: Everything Is W(3,3)}

\section{The Bernoulli Bridge}

The von~Staudt--Clausen theorem selects primes $p$ satisfying $(p{-}1)\mid k$
at weight $k=12$. These are $\{2, 3, 5, 7, 13\}$---exactly the
cyclotomic primes of $\GQ(q,q)$ at $q=3$:
\begin{equation}
\boxed{\den(B_k) = \Phi_1(q)\cdot q\cdot 5\cdot\Phisix\cdot\Phithree = 2730.}
\end{equation}

\begin{theorem}[The Leech Lattice Theta Constant]
\label{thm:leech-theta}
The Leech lattice theta function
$\Theta_{\Lambda_{24}} = E_{12} - \frac{65520}{691}\Delta$
has its constant expressed entirely in W(3,3) terms:
\begin{equation}
\boxed{65520 = f\cdot\den(B_k) = 24\times 2730.}
\end{equation}
\end{theorem}

The Ramanujan discriminant $\Delta = \eta^f$ has exponent $f=24$,
and the denominator $691$ is the irregular-prime numerator of~$B_k$.
The entire formula for the Leech lattice theta function is
\emph{written in W(3,3) parameters}.


\section{The Fine-Structure Constant at $0.2\sigma$}

From prior work in this program (see \texttt{docs/DEEP\_ANALYSIS\_MARCH2026.md}):
\begin{equation}
\alpha^{-1} = (k{-}1)^2 + \mu^2 + \frac{v}{(k{-}1)[(k{-}\lambda)^2{+}1] + \frac{q}{\lambda(k{-}1)}}
= \frac{669969}{4889} = 137.035999182\ldots
\end{equation}
The tree-level value $(k{-}1)^2+\mu^2 = |11{+}4i|^2 = 137$ is the
unique Gaussian integer norm (Fermat two-square theorem), and the
one-loop correction $v/M_{\mathrm{eff}}$ uses the W(3,3) generation
feedback. Agreement with CODATA~2022: $\mathbf{0.2\sigma}$.


\section{The Mass Hierarchy at Exact Match}

The charm-to-top mass ratio is (see \texttt{docs/LINF\_TOWER\_MASS\_DERIVATION.md}):
\begin{equation}
\frac{m_c}{m_t} = \frac{1}{k^2-2\mu} = \frac{1}{136},
\end{equation}
matching $m_c/m_t = 1.27/172.69 = 0.00735$ from PDG to $0.02\%$.
The matrix $G^{136}$ of the $L_\infty$ tower on the W(3,3) chain
complex gives all three quark generation masses.


\section{The Complete Picture}

\renewcommand{\arraystretch}{1.15}
\begin{center}
\begin{tabular}{@{}llll@{}}
\toprule
\textbf{Domain} & \textbf{Object} & \textbf{W(3,3) Formula} & \textbf{Value} \\
\midrule
Number theory & $(q{-}2)!$ & $1$ & $q=3$ \\
Coding theory & Ternary Golay & $[k, q!, q!]_q$ & $[12,6,6]_3$ \\
Coding theory & Binary Golay & $[f, k, 2^q]_2$ & $[24,12,8]_2$ \\
Lattices & Leech kissing & $E q^2\Phisix\Phithree$ & $196560$ \\
Lattices & $\Theta_{\Lambda_{24}}$ const & $f\cdot\den(B_k)/691$ & $65520/691$ \\
Moonshine & $j$ constant & $(2^{q+\lambda}{-}1)f$ & $744$ \\
Moonshine & $c_1(j)$ & $q^2(E\Phisix\Phithree+\mu q^2)$ & $196884$ \\
Modular forms & $\tau(q)$ & $\binom{\Theta}{q+\lambda}$ & $252$ \\
Modular forms & $\Delta$ exponent & $f$ & $24$ \\
Homotopy & $\pi_q^s$ & $\mathbb{Z}/f$ & $\mathbb{Z}/24$ \\
Periodicity & Bott period & $2^q$ & $8$ \\
Algebra & Last NDA dim & $2^q$ & $8$ \\
Topology & Knot dim & $q$ & $3$ \\
Hurwitz & Bound constant & $k\Phisix$ & $84$ \\
Physics & $\alpha^{-1}$ & $|z|^2+v/M_{\mathrm{eff}}$ & $137.036$ \\
Physics & $m_c/m_t$ & $1/(k^2{-}2\mu)$ & $1/136$ \\
Monster & $|M|$ at $3$ & $3^{2\Theta}$ & $3^{20}$ \\
Monster & $|M|$ at $2$ & $2^{v+q!}$ & $2^{46}$ \\
\bottomrule
\end{tabular}
\end{center}

\bigskip
\noindent
\emph{One prime.\ One geometry.\ Everything is arithmetic.}
